\section[A Two-Point Separation Between LOO and Replace-One Stability]{A Two-Point Separation Between LOO and Replace-One Stability}
\label{sec:finite-separation}

We now come to the main explicit witness. The statement is
small, explicit, and completely deterministic.

\begin{note}[Remark on the example]
The example uses no distance tie at the decisive replace-one query; the
separation is therefore not merely an artifact of nearest-neighbor
tie-breaking.
\end{note}

\begin{proposition}[Two-point $1$-NN separation]
\label{prop:6.1}
There exists a finite graph metric space, an ordered labeled sample $S$, and a
query-label pair $(x,y)$ such that
\[
\Delta_{\loo}(S,i)=0 \quad\text{for every sample index } i,
\]
while for some index $j$ and some admissible replacement point $z$,
\[
\Delta_{\rep}(S,j,z,x,y)=1.
\]
\end{proposition}

\begin{note}[Proof sketch]
Let the metric space contain two vertices $a,b$ joined by one edge, so
$d(a,b)=1$. Consider the ordered sample
\[
S=((a,0),(b,0)).
\]
Use deterministic $1$-NN with the fixed tie-breaking conventions from
\cref{sec:definitions}.

For LOO, delete the first occurrence. The remaining sample is $((b,0))$, whose
$1$-NN prediction at the deleted occurrence's location $a$ is still $0$ because
the sole remaining neighbor has label $0$. Hence the loss at $(a,0)$ is
unchanged. The same argument applies when deleting the second occurrence and
evaluating at $(b,0)$. Therefore $\Delta_{\loo}(S,1)=\Delta_{\loo}(S,2)=0$.

Now evaluate at the fixed query-label pair $(x,y)=(a,0)$. Replace the first
occurrence $(a,0)$ by $(a,1)$, obtaining
\[
S^{1\leftarrow (a,1)}=((a,1),(b,0)).
\]
At query $a$, the original sample predicts $0$ because $(a,0)$ is the unique
nearest occurrence. In the perturbed sample, the unique nearest occurrence is
now $(a,1)$, so the prediction becomes $1$. Hence
$\Delta_{\rep}(S,1,(a,1),a,0)=1$.
\end{note}

\Cref{prop:6.1} gives the witness in explicit form. The example is not exotic;
it is simple enough to inspect line-by-line. The phenomenon already appears
before any asymptotics, averages, or large combinatorial constructions enter.

%% Minimality certificate tables moved to supplementary material.
The supplementary material contains a finite-search pipeline that looks for
smaller or structurally simpler witnesses under various restrictions. The
witness in \cref{prop:6.1} is theorem-level because it is explicit and
hand-checkable. Claims that no witness exists below a certain search frontier
are \emph{not} theorem-level; they are computational certificates tied to the
enumerated search space.

This distinction is worth making explicit. A paper can responsibly say ``here is an explicit
witness we know'' and separately say ``here is a reproducible finite
certificate supporting minimality within the searched regime.'' It should not
blur those two statements together.